\section{Design Background}


\subsection{Barske Impeller}

\todo{I just realized I don't have the numbers to back me up that disc friction is better. I will waffle, then fill with literature and numbers later on.}

Centrifugal pumpss have multiple main sources of losses. The one that becomes significant at low specific speeds is the disc friction losses. When a rotor object rotates in a fluid, the shear stresses in the fluid cause a drag torque on the object. In most centrifugal pumps the impeller is disc shaped, and the loss is called disc friction loss. \todo{gulich pdf pg 144 says axial case cleareance only has a very small effect on disc friction. It also says that due to leakage flow, partload recirculation and exchange of momentum, tolerance of calculated disc friction loss can be at least +/-25\%, also, somehow head coefficient is related to disc friction? it says to maximize it for low specific speed}. With low specific speed pumps, disk frction is the main source of loss. At $n_q=10$, $\psi_{\textrm{opt}=1}$ and $\mathrm{Re}=10^7$, the disc friction is approxiomately $30\%$ of the usefulf power of the impeller \cite{gulich page 103}. Gulich suggests that a high head coefficient \todo{high head coeffient not defined yet} reduces the diameter and consequently less disk frciont since $P_{RR}\sim D^5$. This is the reason why low specific speed centrifugal pumps physically cannot achieve the same high efficiencies as higher specific speed pumps. \todo{this paragraph is almost fully copied from gulich, need to rephrase.}

\todo{gulich page 429 says impeller losses are low at low specific speeds, referencing fig 3.33. Oop, fig3.33 shows the disc friction loss i think?}On the other hand, at low specific speeds, impeller losses are low as shown by \cref{fig3.33gulich}. \todo{explain this graph} As a result, surface finish on the impeller blades are largely unimportant, allowing for simpler manufacturing processes. Conversely, collector losses become dominant. \todo{need to explain what a collector is} Gulich explains \cite{gulich 7.12} that for $n_q<12$ even higher efficiencies are expected than with volutes. This simplifies the design and manufacturing of the pump, which is a significant advantage at small scales where tolerances are both harder to achieve but contribute to more significant losses. \todo{need to explain what a volute is} This is important as the high fluid velocity at the peripheri makes the surface finish of the annular casing of primary importance for pump efficiency \cite{lock design features}.

As a result, the Barske impeller which was developed for the need for manufacturing simplicity and low specific speed applications became the optimal choice for the project.